% Generated from locale/tr/content/incompleteness/incompleteness-provability/tarski-thm.tex; do not hand-edit.


\olfileid[tr]{inc}{inp}{tar}

\olsection{Doğruluğun Tanımlanamazlığı}

\emph{Tanımlanabilirlik} kavramı, aritmetik dili için biçimsel bir anlambilime
sahip olmaya dayanır. Aritmetik dilindeki bir formüller ve tümceler kümesini
betimledik. ``Amaçlanan yorum'', böyle tümceleri doğal sayılar hakkında
önermelerde bulunuyormuş gibi okumaktır ve böyle bir önerme doğru ya da yanlış
olabilir. $\Struct{N}$, evreni $\Nat$ olan ve aritmetik dilindeki simgelerin
standart yorumlandığı !!{structure} olsun. O zaman $\Sat{N}{!A}$, ``$!A$
standart yorumda doğrudur'' anlamına gelir.

\begin{defn}
Doğal sayıların bir $R(x_1,\dots,x_k)$ bağıntısı, ancak ve ancak aritmetik
dilinde, her $n_1,\dots,n_k$ için $R(n_1,\dots,n_k)$ ancak ve ancak
$\Sat{N}{!A(\num n_1,\dots,\num n_k)}$ ise geçerli olacak biçimde bir
$!A(x_1,\dots,x_k)$ formülü varsa $\Struct{N}$ içinde \emph{tanımlanabilir}dir.
\end{defn}

Başka türlü söylersek bir bağıntı $\Struct{N}$ içinde ancak ve ancak
$\Th{TA}$ kuramında temsil edilebilirse tanımlanabilirdir; burada
$\Th{TA} = \Setabs{!A}{\Sat{N}{!A}}$, aritmetiğin doğru tümceleri kümesidir.
(Bu size hemen açık gelmiyorsa tanımlara dönüp bunun böyle olduğuna kendinizi
ikna etmelisiniz.)

\begin{lem}
Her hesaplanabilir bağıntı $\Struct{N}$ içinde tanımlanabilirdir.
\end{lem}

\begin{proof}
$\Th{Q}$ içinde bir bağıntıyı temsil eden formülün $\Struct{N}$ içinde aynı
bağıntıyı tanımladığını denetlemek kolaydır.
\end{proof}

Şimdi tersinin de doğru olup olmadığını sorabiliriz. Yani $\Struct{N}$ içinde
tanımlanabilir her bağıntı hesaplanabilir midir? Yanıt hayırdır. Örneğin:

\begin{lem}
Durma bağıntısı $\Struct{N}$ içinde tanımlanabilirdir.
\end{lem}

\begin{proof}
Kleene normal biçim teoreminin, her kısmi hesaplanabilir $f$ fonksiyonu için,
bütün $x \in \Nat$ değerlerinde $f(x)=U(\umin{s}{T(e,x,s)})$ olacak biçimde
bir $e$ indisi bulunduğunu söylediğini anımsayın; burada $U$ ile $T$ ilkel
özyinelemeli ve dolayısıyla totaldir. Öyleyse $f(x)$ ancak ve ancak
(yani hesaplama ancak ve ancak duruyorsa) $T(e,x,s)$'nin geçerli olduğu bir
$s$ varsa tanımlıdır.

Şimdi $H$ durma bağıntısı, yani
\[
H = \Setabs{\tuple{e,x}}{\lexists[s][T(e, x, s)]}.
\]
olsun. $!D_T$, $T$'yi $\Struct{N}$ içinde tanımlasın. O zaman
\[
H = \Setabs{\tuple{e,x}}{\Sat{N}{\lexists[s][!D_T(\num e, \num x, s)]}},
\]
olur; dolayısıyla $\lexists[s][!D_T(z,x,s)]$, $H$'yi $\Struct{N}$ içinde
tanımlar.
\end{proof}

\begin{prob}
$Q(n) \defiff n \in \Setabs{\Gn{!A}}{\Th{Q} \Proves !A}$ bağıntısının
aritmetikte tanımlanabilir olduğunu gösterin.
\end{prob}

Peki $\Th{TA}$'nın kendisi? Aritmetikte tanımlanabilir midir? Başka deyişle
$\Setabs{\Gn{!A}}{\Sat{N}{!A}}$ kümesi aritmetikte tanımlanabilir midir?
Tarski teoremi bu soruyu olumsuz yanıtlar.

\begin{thm}
\ollabel{thm:tarski}
Aritmetiğin doğru !!{sentence}s kümesi aritmetikte tanımlanabilir değildir.
\end{thm}

\begin{proof} 
$!D(x)$'in bu kümeyi tanımladığını, yani $\Sat{N}{!A}$ ancak ve ancak
$\Sat{N}{!D(\gn{!A})}$ ise geçerli olduğunu varsayın. Sabit nokta lemması
uyarınca $\Th{Q} \Proves !A \liff \lnot !D(\gn{!A})$ ve dolayısıyla
$\Sat{N}{!A \liff \lnot !D(\gn{!A})}$ olacak biçimde bir $!A$ formülü vardır.
Fakat o zaman $\Sat{N}{!A}$ ancak ve ancak
$\Sat{N}{\lnot !D(\gn{!A})}$ ise geçerlidir; bu ise $!D(y)$'nin aritmetiğin
doğru önermeleri kümesini tanımladığı varsayımıyla çelişir.
\end{proof}

Tarski bu çözümlemeyi daha genel bir felsefi doğruluk kavramına uyguladı.
Herhangi bir $L$ dili verildiğinde Tarski'ye göre $L$ için yeterli bir doğruluk kavramı, her
$X$ tümcesi için şunu sağlamalıdır:
\begin{quote}
`$X$', ancak ve ancak $X$ ise doğrudur.
\end{quote}
Tarski'nin İngilizce için sıkça alıntılanan örneği şu tümcedir:
\begin{quote}
`Kar beyazdır', ancak ve ancak kar beyazsa doğrudur.
\end{quote}
Ancak köşegen fonksiyonunu temsil edecek kadar güçlü herhangi bir dil ve
herhangi bir dilsel $T(x)$ yüklemi için ``$X$ ancak ve ancak
$T(\text{`$X$'})$ değilse'' koşulunu sağlayan bir $X$ tümcesi kurabiliriz.
Bir doğruluk yükleminin bazı tümceleri hem doğru hem yanlış ilan etmesini
istemediğimizden Tarski, bir dildeki bütün tümceler için doğruluk yükleminin,
bir biçimde dilin sınırlarının dışına çıkmadan belirtilemeyeceği sonucuna
vardı. Başka deyişle bir dilin doğruluk yüklemi o dilin kendisinde
tanımlanamaz.

